\documentclass[10pt,a4paper]{amsart}
\usepackage[margin=19mm]{geometry}
\usepackage{amsmath,amssymb,mathtools}
\usepackage[scr=rsfs,cal=euler]{mathalfa}
\usepackage{xfrac}
\usepackage[unicode,hidelinks,bookmarks=false]{hyperref}
\newcommand{\ie}{i.e.\ }
\newcommand{\cf}{cf.\ }
\newcommand{\resp}{respectively\ }
\newcommand{\Subsection}{\subsection}
\providecommand{\nobreakdash}{\nobreak\hbox{-}}
\setlength{\emergencystretch}{2em}
\multlinegap0pt
\usepackage{mathtools}
\usepackage{enumitem}
\usepackage{cleveref}
\hypersetup{
  colorlinks=true,
  linkcolor=black,
  citecolor=blue,
  urlcolor=blue
}
\newcommand{\ii}{\mathrm{i}}
\newcommand{\dd}{\mathrm{d}}
\newcommand{\R}{\mathbb{R}}
\newcommand{\N}{\mathbb{N}}
\newcommand{\cR}{\mathcal{R}}
\newcommand{\cP}{\mathcal{P}}
\newcommand{\half}{\frac12}
\newcommand{\floor}[1]{\left\lfloor #1\right\rfloor}
\newcommand{\poch}[2]{\left(#1\right)_{#2}}
\numberwithin{equation}{section}
\newtheorem{theorem}{Theorem}[section]
\newtheorem{proposition}[theorem]{Proposition}
\newtheorem{corollary}[theorem]{Corollary}
\newtheorem{remark}[theorem]{Remark}
\begin{document}
\title[Morse momentum wavefunctions and rational functions]{Morse momentum wavefunctions and rational functions}
\author{Luc Vinet, Alexei Zhedanov}
\date{}
\maketitle
\noindent\emph{Source and licence.} Morse momentum wavefunctions and rational functions. Version arXiv:2606.24401v1 (CC BY 4.0). This material is distributed under the Creative Commons Attribution 4.0 International licence, http://creativecommons.org/licenses/by/4.0/, as recorded in the arXiv OAI licence metadata for this exact version and shown on the abstract page https://arxiv.org/abs/2606.24401v1. Full rights notice: licenses/arxiv-2606.24401-CC-BY-4.0.txt.\par\smallskip
\noindent\textit{Editorial scope.} Counted unit: the original \texttt{theorem} environment \texttt{thm:zeros} at source lines 554--561 together with its complete proof at lines 563--571; the delivered document keeps source lines 154--572 so that every referenced label and every citation resolves inside it. Native editable source is the paper's own concatenated TeX; the AMS wrapper is editorial formatting only. No publisher class, upstream script or unrelated artwork is redistributed.
\begin{equation}
  E_n=-\left(g-n-\frac12\right)^2,
  \qquad n=0,1,\dots,N,
  \qquad N=\floor{g-\frac12}.
  \label{eq:energies}
\end{equation}
The corresponding coordinate wavefunctions can be written as
\begin{equation}
  \psi_n(x)=G_n e^{-y/2} y^{g-n-1/2}
  {}_1F_1\!\left(
  \begin{matrix}-n\\ 2g-2n\end{matrix};y
  \right),
  \qquad y=2g e^{-x},
  \label{eq:coordinate_wavefunctions}
\end{equation}
where $G_n$ is a normalization constant.

The momentum representation is defined by the Fourier transform
\begin{equation}
  \Psi(p)=\frac{1}{\sqrt{2\pi}}
  \int_{-\infty}^{\infty} e^{-\ii px}\psi(x)\,\dd x.
  \label{eq:fourier}
\end{equation}
In momentum space the coordinate operator becomes
\begin{equation}
  x=\ii\partial_p.
\end{equation}
Consequently the exponentials in the Morse potential act as shift operators,
\begin{equation}
  e^{-x}=T,
  \qquad e^{-2x}=T^2,
  \label{eq:shift_exp}
\end{equation}
where
\begin{equation}
  T^jF(p)=F(p-j\ii),
  \qquad j=0,\pm1,\pm2,\dots .
  \label{eq:shift_def}
\end{equation}
Thus the Hamiltonian becomes
\begin{equation}
  H=p^2+g^2(T^2-2T),
  \label{eq:momentum_hamiltonian}
\end{equation}
and the bound-state equation reads
\begin{equation}
  p^2\Psi_n(p)+g^2\bigl(\Psi_n(p-2\ii)-2\Psi_n(p-\ii)\bigr)
  =E_n\Psi_n(p).
  \label{eq:momentum_difference_eq}
\end{equation}
The explicit expression obtained in \cite{SunDong2012} is
\begin{equation}
  \Psi_n(p)=K_n\Gamma\!\left(g-n-\frac12+\ii p\right)g^{-\ii p}
  {}_2F_1\!\left(
  \begin{matrix}
  -n,\; g-n-\frac12+\ii p\\
  2g-2n
  \end{matrix};2
  \right),
  \label{eq:sun_dong}
\end{equation}
where $K_n$ is independent of $p$.

\section{Rational eigenfunctions}\label{sec:rational}

The ground-state momentum wavefunction is, up to normalization,
\begin{equation}
  \Psi_0(p)=g^{-\ii p}\Gamma\!\left(g-\frac12+\ii p\right).
  \label{eq:ground_state}
\end{equation}
It is a solution of \eqref{eq:momentum_difference_eq} with energy
\begin{equation}
  E_0=-\left(g-\frac12\right)^2.
\end{equation}
We write the excited states in the form
\begin{equation}
  \Psi_n(p)=\Psi_0(p)R_n(p).
  \label{eq:factorization}
\end{equation}
Substitution in \eqref{eq:momentum_difference_eq} gives the eigenvalue equation
\begin{equation}
  W R_n(p)=E_n R_n(p),
  \label{eq:W_eigenvalue}
\end{equation}
where the difference operator $W$ is
\begin{equation}
  W=p^2I+\bigl((g+\ii p)^2-\frac14\bigr)T^2
  -2g\left(g-\frac12+\ii p\right)T.
  \label{eq:W_operator}
\end{equation}
In particular,
\begin{equation}
  W\{1\}=E_0.
\end{equation}

Let
\begin{equation}
  \chi_k(p)=\frac{1}{\left(-\ii p-g+\frac32\right)_k},
  \qquad k=0,1,2,\dots,
  \label{eq:chi_basis}
\end{equation}
where
\begin{equation}
  (a)_k=a(a+1)\cdots(a+k-1),
  \qquad (a)_0=1.
\end{equation}
For each $n$, the functions $\chi_0,\chi_1,\dots,\chi_n$ span the finite-dimensional space
\begin{equation}
  \cR_n=\operatorname{span}\{\chi_0(p),\chi_1(p),\dots,\chi_n(p)\}.
\end{equation}
Equivalently, elements of $\cR_n$ are rational functions of the form
\begin{equation}
  \frac{Q_n(p)}{
  \left(-\ii p-g+\frac32\right)
  \left(-\ii p-g+\frac52\right)
  \cdots
  \left(-\ii p-g+\frac{2n+1}{2}\right)},
  \label{eq:rational_flag}
\end{equation}
where $Q_n(p)$ is a polynomial of degree at most $n$.

\begin{proposition}\label{prop:bidiagonal}
The operator $W$ acts on the basis \eqref{eq:chi_basis} as
\begin{equation}
  W\chi_k(p)=
  -\left(g-k-\frac12\right)^2\chi_k(p)+2k\chi_{k-1}(p),
  \qquad k\ge 0,
  \label{eq:bidiagonal_action}
\end{equation}
with the convention $\chi_{-1}=0$.
\end{proposition}

\begin{proof}
This follows by direct substitution.  Under the shift $T$, the variable $\ii p$ is replaced by
$\ii p+1$, and hence
\begin{equation}
  T\chi_k(p)=\frac{1}{\left(-\ii p-g+\frac12\right)_k},
  \qquad
  T^2\chi_k(p)=\frac{1}{\left(-\ii p-g-\frac12\right)_k}.
\end{equation}
Substituting these expressions in \eqref{eq:W_operator} and reducing the resulting rational
function to the basis $\{\chi_k,\chi_{k-1}\}$ gives \eqref{eq:bidiagonal_action}.
\end{proof}

It follows that, for each $n=0,1,\dots,N$, there is an eigenfunction of $W$ in $\cR_n$ with
eigenvalue $E_n$.  Write
\begin{equation}
  R_n(p)=\sum_{k=0}^{n} A_{n,k}\chi_k(p).
  \label{eq:R_expansion}
\end{equation}
Using \eqref{eq:bidiagonal_action} in \eqref{eq:W_eigenvalue} gives the recurrence
\begin{equation}
  A_{n,k+1}
  =\frac{(k-n)(k+n+1-2g)}{2(k+1)}A_{n,k},
  \qquad k=0,1,\dots,n-1.
  \label{eq:A_recurrence}
\end{equation}
Choosing $A_{n,0}=1$, one obtains
\begin{equation}
  A_{n,k}=\frac{(-n)_k(n+1-2g)_k}{k!}\left(\frac12\right)^k.
  \label{eq:A_coefficients}
\end{equation}
Hence
\begin{equation}
  R_n(p)=
  {}_2F_1\!\left(
  \begin{matrix}
  -n,\; n+1-2g\\
  -\ii p-g+\frac32
  \end{matrix};\frac12
  \right).
  \label{eq:R_hypergeometric}
\end{equation}
Thus
\begin{equation}
  \Psi_n(p)=\kappa_n\Psi_0(p)R_n(p)
  \label{eq:Psi_R_solution}
\end{equation}
provides the bound-state solutions of \eqref{eq:momentum_difference_eq}, with energies
\eqref{eq:energies}.

\section{Identification with finite biorthogonal rational functions}\label{sec:kmj}

The rational functions \eqref{eq:R_hypergeometric} are naturally related to the finite family of
rational functions introduced by Koepf and Masjed-Jamei \cite{KoepfMasjedJamei2007}.  In the
notation used here, that family is
\begin{equation}
  B_n(x;a,b)=
  {}_2F_1\!\left(
  \begin{matrix}
  -n,\; n-a-b\\
  1-a-x
  \end{matrix};\frac12
  \right).
  \label{eq:KMJ_functions}
\end{equation}
The functions $B_n$ satisfy a finite biorthogonality relation of the form
\begin{equation}
  \int_{-\infty}^{\infty}
  \Gamma(a+\ii x)\Gamma(b-\ii x)
  B_n(\ii x;a,b)B_m(-\ii x;b,a)\,\dd x=0,
  \qquad n\ne m,
  \label{eq:KMJ_biorthogonality}
\end{equation}
for the admissible finite range of indices.  In the symmetric specialization $a=b$, this
biorthogonality becomes an ordinary orthogonality relation.

Comparing \eqref{eq:R_hypergeometric} and \eqref{eq:KMJ_functions}, one obtains
\begin{equation}
  x=\ii p,
  \qquad a=b=g-\frac12.
  \label{eq:parameter_identification}
\end{equation}
Equivalently,
\begin{equation}
  R_n(p)=B_n\!\left(\ii p;g-\frac12,g-\frac12\right).
  \label{eq:R_B_identification}
\end{equation}

\begin{theorem}\label{thm:main_identification}
The rational factors of the momentum-space bound states of the Morse potential coincide with the
symmetric Koepf--Masjed-Jamei finite rational functions:
\begin{equation}
  R_n(p)=B_n\!\left(\ii p;g-\frac12,g-\frac12\right),
  \qquad n=0,1,\dots,N.
\end{equation}
Consequently, the Morse momentum wavefunctions provide a physical realization of a finite family
of orthogonal rational functions belonging to the $R_{II}$ framework.
\end{theorem}

The physical orthogonality of the Morse bound states gives the corresponding orthogonality
relation for these rational functions.  With the normalization chosen so that the wavefunctions are
orthonormal,
\begin{equation}
  \int_{-\infty}^{\infty}
  \overline{\Psi_m(p)}\Psi_n(p)\,\dd p=\delta_{mn}.
\end{equation}
Using \eqref{eq:factorization}, this becomes
\begin{equation}
  \int_{-\infty}^{\infty}
  w(p)R_n(p)\overline{R_m(p)}\,\dd p=h_n\delta_{mn},
  \label{eq:R_orthogonality_physical}
\end{equation}
where
\begin{equation}
  w(p)=|\Psi_0(p)|^2
  =\left|\Gamma\!\left(g-\frac12+\ii p\right)\right|^2.
  \label{eq:weight}
\end{equation}
Here the factor $g^{-\ii p}$ has modulus one for real $p$ and $g>0$.  Since
$\overline{R_m(p)}=R_m(-p)$ for real $p$, \eqref{eq:R_orthogonality_physical} is the symmetric
specialization of the Koepf--Masjed-Jamei biorthogonality.

\begin{remark}
The terminology $R_{II}$ is often attached to rational functions whose recurrence relations are
encoded by tridiagonal generalized eigenvalue problems.  In its standard form, an $R_{II}$
recurrence contains a quadratic factor multiplying the preceding element.  Equivalently, it can be
written as a tridiagonal matrix pencil.  The Koepf--Masjed-Jamei functions form a finite
hypergeometric realization of this rational framework; the Morse model realizes the symmetric
orthogonal specialization of that family.
\end{remark}

\section{Generalized eigenvalue problem and bispectrality}\label{sec:gevp}

The rational functions $R_n(p)$ satisfy two eigenvalue problems.  The first is the difference
equation \eqref{eq:W_eigenvalue}, where the operator $W$ acts on the momentum variable $p$ and
the eigenvalue is the energy $E_n$.

The second is a generalized eigenvalue problem in the degree variable.  This is the rational
analogue of the three-term recurrence relation for orthogonal polynomials.  In the theory of
biorthogonal rational functions, such recurrences are naturally written as matrix pencils
\begin{equation}
  J_1 R=zJ_2 R,
  \label{eq:abstract_GEVP}
\end{equation}
where $J_1$ and $J_2$ are tridiagonal matrices \cite{Zhedanov1999}.  For the Morse rational
functions, the spectral parameter is $z=\ii p$ and the dual problem is
\begin{equation}
  J_1R(p)=\ii p\,J_2R(p),
  \label{eq:morse_GEVP}
\end{equation}
where
\begin{equation}
  R(p)=\bigl(R_0(p),R_1(p),\dots,R_N(p)\bigr)^T.
\end{equation}
In components, for the interior values of the finite system,
\begin{align}
  &\eta_{1,n}R_{n+1}(p)+\eta_{2,n}R_{n-1}(p)+\eta_{0,n}R_n(p)
  \notag\\
  &\qquad
  =-\ii p\left(
  \xi_{1,n}R_{n+1}(p)+\xi_{2,n}R_{n-1}(p)+\xi_{0,n}R_n(p)
  \right).
  \label{eq:component_GEVP}
\end{align}
The coefficients are obtained from contiguous relations for the Gauss hypergeometric function:
\begin{align}
  \eta_{2,n}&=n-g-\frac12,
  &
  \eta_{1,n}&=
  \frac{(g-n)(2g-n-1)(g-n-\frac32)}{n(g-n-1)},
  &
  \eta_{0,n}&=
  \frac{g(g-n-\frac12)}{n(g-n-1)},
  \label{eq:eta_coefficients}
\end{align}
while
\begin{align}
  \xi_{2,n}&=1,
  &
  \xi_{1,n}&=
  \frac{(g-n)(2g-n-1)}{n(g-n-1)},
  &
  \xi_{0,n}&=-
  \frac{g(2g-2n-1)}{n(g-n-1)}.
  \label{eq:xi_coefficients}
\end{align}
Boundary terms are obtained by truncating the finite system at $n=0$ and $n=N$.

Equations \eqref{eq:W_eigenvalue} and \eqref{eq:morse_GEVP} exhibit the bispectral character of
the Morse rational functions.  The difference operator $W$ acts in the momentum variable and has
spectrum $E_n$, whereas the generalized eigenvalue problem acts in the degree variable and has
spectral parameter $\ii p$.  This is the rational counterpart of the bispectrality of classical
orthogonal polynomials, with tridiagonal matrix pencils replacing ordinary three-term recurrence
relations.

\section{Relation with Meixner--Pollaczek polynomials and zeros}\label{sec:mp}

There is another useful representation of the momentum wavefunctions.  We use the transformation
formula
\begin{equation}
  {}_2F_1\!\left(
  \begin{matrix}
  -n,\;\beta\\
  \gamma
  \end{matrix};z
  \right)
  =\frac{(\beta)_n}{(\gamma)_n}(-z)^n
  {}_2F_1\!\left(
  \begin{matrix}
  -n,\;1-\gamma-n\\
  1-\beta-n
  \end{matrix};z^{-1}
  \right).
  \label{eq:hypergeometric_transform}
\end{equation}
Applying this transformation to \eqref{eq:R_hypergeometric} gives
\begin{equation}
  R_n(p)=
  \frac{(-1/2)^n(n+1-2g)_n}{(3/2-g-\ii p)_n}
  V_n(p),
  \label{eq:R_V_factorization}
\end{equation}
where
\begin{equation}
  V_n(p)=
  {}_2F_1\!\left(
  \begin{matrix}
  -n,\; g-n+\ii p-\frac12\\
  2g-2n
  \end{matrix};2
  \right).
  \label{eq:V_def}
\end{equation}
The function $V_n(p)$ is a polynomial in $p$ of degree $n$.  In this form all poles of
$R_n(p)$ are contained in the factor $(3/2-g-\ii p)_n^{-1}$, and \eqref{eq:V_def} agrees with the
hypergeometric expression for the momentum wavefunction displayed in \eqref{eq:sun_dong}.

The Meixner--Pollaczek polynomials are defined, up to a normalization factor independent of the
argument, by
\begin{equation}
  P_n(t;\mu,\phi)=
  {}_2F_1\!\left(
  \begin{matrix}
  -n,\;\mu+\ii t\\
  2\mu
  \end{matrix};1-e^{-2\ii\phi}
  \right).
  \label{eq:MP_def}
\end{equation}
Comparison of \eqref{eq:V_def} with \eqref{eq:MP_def} gives
\begin{equation}
  \phi=\frac{\pi}{2},
  \qquad \mu=g-n,
  \qquad t=p+\frac{\ii}{2}.
  \label{eq:MP_parameters}
\end{equation}
Therefore
\begin{equation}
  V_n(p)=P_n\!\left(p+\frac{\ii}{2};g-n,\frac{\pi}{2}\right),
  \label{eq:V_MP}
\end{equation}
up to the normalization convention used for the Meixner--Pollaczek polynomials.

For $\mu>0$, the Meixner--Pollaczek polynomials are positive definite and have $n$ simple real
zeros.  In the present case $\mu=g-n$, and for bound states
$n=0,1,\dots,\floor{g-1/2}$ one has $g-n>0$.
Consequently, the zeros of $V_n(p)$ are obtained from the real zeros of a positive definite
Meixner--Pollaczek polynomial by the shift $t=p+\ii/2$.


\begin{theorem}\label{thm:zeros}
For each bound state $n=0,1,\dots,N$, the momentum-space wavefunction $\Psi_n(p)$ has exactly
$n$ simple zeros.  They lie on the horizontal line
\begin{equation}
  \operatorname{Im}p=-\frac12
\end{equation}
in the complex $p$-plane.  In particular, $\Psi_n(p)$ has no zeros on the real momentum axis.
\end{theorem}

\begin{proof}
The ground-state factor \eqref{eq:ground_state} has no zeros, and the factor
$(3/2-g-\ii p)_n^{-1}$ in \eqref{eq:R_V_factorization} has no zeros.  Hence the zeros of
$\Psi_n(p)$ are precisely the zeros of $V_n(p)$.  By \eqref{eq:V_MP}, these are obtained from the
real zeros of the Meixner--Pollaczek polynomial $P_n(t;g-n,\pi/2)$ by the shift $t=p+\ii/2$.
Thus real zeros in $t$ correspond to zeros of $\Psi_n(p)$ lying on
$\operatorname{Im}p=-1/2$.  Simplicity follows from positive definiteness of the
Meixner--Pollaczek orthogonality for $g-n>0$.
\end{proof}


\begin{thebibliography}{99}
\bibitem[KoepfMasjedJamei2007]{KoepfMasjedJamei2007} W. Koepf and M. Masjed-Jamei, Two classes of special functions using Fourier transforms of some finite classes of classical orthogonal polynomials, \emph{Proc. Amer. Math. Soc.} \textbf{135} (2007), 3599--3606.
\bibitem[SunDong2012]{SunDong2012} G. H. Sun and S. H. Dong, Morse potential in the momentum representation, \emph{Commun. Theor. Phys.} \textbf{58} (2012), 815--818.
\bibitem[Zhedanov1999]{Zhedanov1999} A. Zhedanov, Biorthogonal rational functions and the generalized eigenvalue problem, \emph{J. Approx. Theory} \textbf{101} (1999), 303--329.
\end{thebibliography}
\end{document}
